%\documentclass[12pt,oneside]{amsart}
\documentclass{scrartcl}

%\documentclass{article}

\usepackage{amsthm,amsmath,amssymb}
\usepackage[numbers]{natbib}

%\usepackage[colorlinks]{hyperref}
\usepackage{hyperref}
\hypersetup{
    colorlinks,%
    citecolor=black,%
    filecolor=black,%
    linkcolor=black,%
    urlcolor=black
}
\usepackage{hypernat}
\usepackage[top=2.5cm, bottom=2.5cm, left=2.8cm,
right=2.8cm]{geometry}
\usepackage{graphicx}


%\setlength{\parskip}{0pt}
%\setlength{\parsep}{0pt}
%\setlength{\headsep}{0pt}
%\setlength{\topskip}{0pt}
%\setlength{\topmargin}{0pt}
%\setlength{\topsep}{0pt}
%\setlength{\partopsep}{0pt}

\linespread{1}

%\usepackage{marvosym}

%\textwidth = 6.5 in \textheight = 9.2 in \oddsidemargin = 0.0 in
%\evensidemargin = 0.0 in \topmargin = -0.0 in \headheight = 0.0 in
%\headsep = 0.0in \parskip = 0.0in \parindent = 0.0in
%\def\baselinestretch{0.871}


\newtheorem{theorem}{Theorem}[section]
\newtheorem{proposition}[theorem]{Proposition}
\newtheorem{problem}[theorem]{Problem}
\newtheorem{fact}[theorem]{Fact}
\newtheorem{lemma}[theorem]{Lemma}
\newtheorem{defn}[theorem]{Definition}
\newtheorem{corollary}[theorem]{Corollary}
\newtheorem{remark}{Remark}[section]
\newtheorem{example}[theorem]{Example}

\newcommand{\E}{{\mathbb E}}
\newcommand{\se}{{\mathcal{E}}}
\newcommand{\R}{{\mathbb R}}
\renewcommand{\P}{{\mathbb P}}
\newcommand{\ac}{{\mathcal{A}}}
\newcommand{\A}{{\mathcal{A}}}
\newcommand{\B}{{\mathcal{B}}}
\newcommand{\C}{{\mathcal{C}}}
\newcommand{\M}{{\mathcal{M}}}
\newcommand{\Mf}{{\mathfrak{M}}}
\newcommand{\Qcal}{{\mathcal{Q}}}
\newcommand{\T}{{\mathcal{T}}}
\newcommand{\Z}{{\mathcal{Z}}}
\newcommand{\K}{{\mathcal{K}}}
\newcommand{\lc}{{\mathcal{L}}}
\newcommand{\Ps}{{\mathcal{P}}}
\newcommand{\G}{{\mathcal{G}}}
\newcommand{\pa}{{\mathring{p}}}
\newcommand{\F}{{\cal F}}
\newcommand{\samp}{{\mathcal{X}}}
\newcommand{\X}{{\mathcal{X}}}
\newcommand{\hi}{{\hat{\Phi}}}
\newcommand{\data}{{\text{data}}}
\newcommand{\relu}{{\text{ReLU}}}

\newcommand{\ridge}{{\hat{\beta}_{\text{ridge}}(\lambda)}}
\newcommand{\lasso}{{\hat{\beta}_{\text{lasso}}(\lambda)}}
\newcommand{\ridgemu}{{\hat{\mu}_t^{\text{ridge}}(\lambda)}}
\newcommand{\lassomu}{{\hat{\mu}_t^{\text{lasso}}(\lambda)}}


\newcounter{rcnt}[section]
\renewcommand{\thercnt}{(\roman{rcnt})}

\def\qt#1{\qquad\text{#1}}

\def\argmin{\mathop{\rm argmin}}
\def\argmax{\mathop{\rm argmax}}


\setlength{\parskip}{1.4 \medskipamount}
%\sloppy
%\linespread{1.3}

\begin{document}

\title{STAT 238 - Bayesian Statistics  \\
  Lecture Thirty Eight} 
\subtitle{Spring 2026, UC Berkeley}

\author{Aditya Guntuboyina}


\date{01 May 2026}

\maketitle

As an example of variational inference, we shall look at an example of
Variational Auto-Encoders (VAE) which are used for generative modeling.

\section{Variational AutoEncoders}

Consider data points $y_1, \dots, y_n$, where each $y_i$ is a binary
vector in $\{0, 1\}^{d_y}$. One can think of each $y_i$ as an image,
flattened from its 2D pixel grid into a single vector of length $d_y$,
where each entry corresponds to one pixel. In general, pixel values
can take many forms---for instance, grayscale images use intensities
in $\{0, 1, \dots, 255\}$, and color images use three such values per
pixel (one per RGB channel). Here we focus on the simplest case:
\emph{binary} images (also called black-and-white or bilevel images),
in which each pixel is either off ($0$) or on ($1$). For example, a
$28 \times 28$ binary image yields $d_y = 28^2 = 784$.

We build up a model for the data in stages. Since each coordinate of
$y_i$ is binary, the most basic assumption is
\begin{align*}
  y_i \mid p_i \overset{\text{independent}}{\sim} \mathrm{Bernoulli}(p_i),
\end{align*}
where $p_i \in [0, 1]^{d_y}$ specifies the probability that each pixel
is on (Bernoulli is interpreted componentwise, so pixels are
conditionally independent given $p_i$). In other words, we are
assuming that each $y_i \in \{0, 1\}^{d_y}$ is Bernoulli with its own
parameter $p_i \in [0, 1]^{d_y}$. We need to further model $p_1,
\dots, p_n$ (otherwise, there will be too many parameters to yield
anything useful). Here it is more convenient to work
on the logit scale: let 
\begin{align*}
  \eta_i = \log\!\frac{p_i}{1 - p_i} \in \mathbb{R}^{d_y},
\end{align*}
applied componentwise.

We shall assume $\eta_1, \dots, \eta_n$ are i.i.d with some common
distribution on $\R^{d_y}$. Further, we assume that this common
distribution, despite living in $\mathbb{R}^{d_y}$, is in fact
concentrated on (or near) a much lower-dimensional set. A
natural way to encode this is to write 
\begin{align*}
  \eta_i = f_\theta(z_i) ~~ \text{ where } z_i
  \overset{\text{i.i.d.}}{\sim} N(0, I_d), 
\end{align*}
where $f_\theta : \mathbb{R}^d \to \mathbb{R}^{d_y}$ is a neural
network with parameters $\theta$ and $d \ll d_y$. The pushforward of
$N(0, I_d)$ under $f_\theta$ defines a $d$-dimensional family of
distributions on $\mathbb{R}^{d_y}$, and the flexibility of $f_\theta$
ensures this family is rich enough to model the structure we care
about. The Gaussian prior on $z_i$ is essentially a convention: any
continuous distribution on $\mathbb{R}^d$ can be written as a
deterministic transformation of a standard Gaussian, so fixing $z_i
\sim N(0, I_d)$ and letting $f_\theta$ be flexible loses no generality
(it is also analytically convenient, as we will see when designing an
inference procedure). Putting the pieces together, and writing
$\sigma$ for the componentwise sigmoid, the overall model can be
written in the following alternative way:  
\begin{align*}
  z_i \overset{\text{i.i.d}}{\sim} N(0, I_d) ~~ \text{ and } ~~ 
  y_i \mid z_i \sim
  \mathrm{Bernoulli}\bigl(\sigma(f_\theta(z_i))\bigr)
\end{align*}
with, for example,
\begin{align}\label{fth_nn}
f_\theta(z_i) = B~\text{ReLU}(A z_i + a) + b ~~ \text{ with }
\theta = (A, a, B, b)
\end{align}
The unknown parameters in this model are the neural network
parameters $\theta$. The latent variables $z_1, \dots, z_n$ are also
unobserved. The $z_i$'s can be integrated out and the model becomes
\begin{align*}
  y_i \overset{\text{i.i.d}}{\sim} p_{\theta}
\end{align*}
where
\begin{align*}
  p_{\theta}(y) = \int p_{\theta}(y \mid z) \varphi(z) dz
\end{align*}
with $\varphi(z)$ being the density of $N(0, I_d)$, and
\begin{align*}
  p_{\theta}(y \mid z) = \prod_{j=1}^{d_y} \left(\sigma(f_{\theta,
  j}(z)) \right)^{y_j} 
  \left(1 - \sigma(f_{\theta,
  j}(z))\right)^{1 - y_j}  ~~ \text{ for } ~~ y = (y_1, \dots,
  y_{d_y}) \in \{0, 
  1\}^{d_y}. 
\end{align*}
The likelihood then is given by: 
\begin{align*}
 f_{\text{data} \mid \theta}(y_1, \dots, y_n) = \prod_{i=1}^n \left[
  \int p_{\theta}(y_i \mid z_i) \varphi(z_i) dz_i \right] 
\end{align*}
and the corresponding log-likelihood is
\begin{align}\label{loglik_actual}
 \log f_{\text{data} \mid \theta}(y_1, \dots, y_n) = \sum_{i=1}^n \log \left[ \int p_{\theta}(y_i \mid z_i) \varphi(z_i)
  dz_i \right].  
\end{align}
This log-likelihood cannot be directly used as the objective function
in optimization software such as PyTorch because of the presence of
the integral. A natural idea is to discretize the integral by sampling
$z^{(l)}_i, l = 1, \dots, M$  from $N(0, I_d)$ and forming the
approximate log-likelihood:
\begin{align}\label{loglik_approx1}
 \log f_{\text{data} \mid \theta}(y_1, \dots, y_n) \approx
  \sum_{i=1}^n \log \left[ \frac{1}{M} \sum_{l = 1}^M  
  p_{\theta}(y_i \mid z_i^{(l)}) \right]. 
\end{align}
The above term \eqref{loglik_approx1} is unlikely to be a good
approximation to the actual log-likelihood \eqref{loglik_actual}. This
 is because the integral $\int p_{\theta}(y_i \mid z_i) \varphi(z_i) dz_i$
in \eqref{loglik_actual} is dominated by the region where
$p_{\theta}(y_i \mid z_i)$ is large i.e., by the posterior
$p_{\theta}(z_i \mid y_i)$. This posterior typically is very
concentrated in the sense that only a tiny region of $z_i \in \R^d$
yields non-negligibe values of $p_{\theta}(y_i \mid
z_i)$. The prior $z_i \sim N(0, I_d)$, by contrast, spreads mass
diffusely over $\R^d$. The consequence will be that almost every
sample $z_i^{(l)}$ lands in a region where $p_{\theta}(y_i \mid
z_i^{(l)})$ is small, so that the average in \eqref{loglik_approx1} is
dominated by one or two (or even zero) lucky samples. This makes the
average in \eqref{loglik_approx1} a highly variance estimate of the
integral in \eqref{loglik_actual}, and hence unreliable as an
optimization objective.

A better approach is to use ideas from variational inference and to
maximize the ELBO instead.

\section{Use of the ELBO for maximization of $\log f_{\text{data} \mid
    \theta}(y_1, \dots, y_n)$ } 

We shall use the following formula whose proof we saw in Lecture 36
(see also Lecture 37): 
\begin{align}\label{evid_elbo}
  \log f_{\text{data} \mid \theta}(y_1, \dots, y_n) = \max_{q}
  \text{ELBO}(q, \theta)
\end{align}
where
\begin{align*}
  & \text{ELBO}(q, \theta) \\ &= \int \dots \int q(z_1, \dots, z_n) \log
  \frac{\prod_{i=1}^np_{\theta}(y_i,  z_i)}{q(z_1, \dots, z_n)}
                           dz_1 \dots dz_n \\
  &= \int \dots \int q(z_1, \dots, z_n) \log \prod_{i=1}^n
    p_{\theta}(y_i \mid z_i) dz_1 \dots dz_n - \int \dots \int q(z_1,
    \dots, z_n) \log \frac{q(z_1, \dots, z_n)}{\varphi(z_1), \dots,
    \varphi(z_n)} dz_1 \dots dz_n \\
  &= \sum_{i=1}^n \int \log p_{\theta}(y_i \mid z_i) q_i(z_i) dz_i -
    KL(q \| \varphi \times \dots \times \varphi).
\end{align*}
where $q_i$ is the $i$-th marginal corresponding to $q$, and $\varphi
\times \dots \times \varphi$ denotes the distribution of $(z_1, \dots,
z_n)$ corresponding to $z_i \overset{\text{i.i.d}}{\sim} N(0, I_d)$.

The maximization in \eqref{evid_elbo} is over all probability
densities $q$ of $z_1, \dots, z_n$. We have also seen (again in Lecture 36)
that the maximum in \eqref{evid_elbo} is achieved when $q$ is the
conditional density of $z_1, \dots, z_n$  given $y_1, \dots, y_n$ and
$\theta$:
\begin{align*}
  q^*(z_1, \dots, z_n) \propto \prod_{i=1}^n\left[ p_{\theta}(y_i \mid z_i)
  \varphi(z_i) \right]. 
\end{align*}
It is clear that this $q^*$ is a product measure $\prod_{i=1}^n
q_i^*(z_i \mid y_i)$ where
\begin{align*}
 q_i^*(z_i \mid y_i)  \propto p_{\theta}(y_i \mid z_i)
  \varphi(z_i) \implies  q_i^*(z_i \mid y_i)  = \frac{p_{\theta}(y_i \mid z_i)
  \varphi(z_i)}{\int p_{\theta}(y_i \mid z_i)
  \varphi(z_i) dz_i}. 
\end{align*}
However because of the somewhat complicated form of $p_{\theta}(y_i
\mid z_i)$, the density $q_i^*$ above is not of a simple form (for
example, it is not a Gaussian). We therefore take a simpler class of
densities $\Qcal$ and then maximize the ELBO over
$\Qcal$. Specifically, we shall take $\Qcal$ to be the class of all
product densities $\prod q_{i, \phi}(z_i \mid y_i)$ for which each
marginal $q_{\phi}(z_i \mid y_i)$ is Gaussian: 
\begin{align*}
  q_{\phi}(z_i \mid y_i) = ~\text{density of}~ N(\mu_{\phi}(y_i),
  \Sigma_{\phi}(y_i)) 
\end{align*}
with
\begin{align*}
  \mu_{\phi}(y_i) = (\mu_{\phi, j}(y_i), j = 1, \dots, d) ~~ \text{
  and } ~~ \Sigma_{\phi}(y_i) = \text{diag} \left(\sigma^2_{\phi, j}(y_i),
  j = 1, \dots, d \right). 
\end{align*}
With this choice $\Qcal$, the ELBO becomes
\begin{align*}
  \text{ELBO}(q, \theta) = \sum_{i=1}^n \int \log p_{\theta}(y_i \mid
  z_i) q_{\phi}(z_i \mid y_i) dz_i - \sum_{i=1}^n KL(q(\cdot \mid y_i)
  \| \phi),
\end{align*}
where we use the fact that the KL between two product distributions
becomes the sum of the KL between the marginals. We can further write
this as: 
\begin{align*}
  \text{ELBO}(q, \theta) = \sum_{i=1}^n \E_{z_i \sim q_{\phi}(\cdot
  \mid y_i)} \log p_{\theta}(y_i \mid z_i) - \sum_{i=1}^n
  KL(N(\mu_{\phi}(y_i), \Sigma_{\phi}(y_i))
  \| N(0, I_d)). 
\end{align*}
The second term above can be written in closed form as (see e.g.,
\url{https://en.wikipedia.org/wiki/Kullback%E2%80%93Leibler_divergence#Examples})
\begin{align*}
  KL(N(\mu_{\phi}(y_i), \Sigma_{\phi}(y_i))
  \| N(0, I_d)) = \frac{1}{2} \sum_{j=1}^d \left\{\sigma^2_{\phi,
  j}(y_i) + \mu^2_{\phi, j}(y_i) - \log \sigma^2_{\phi, j}(y_i) - 1 \right\}.
\end{align*}
For the first term in the ELBO expression above, we need to compute
\begin{align*}
  \E_{z_i \sim q_{\phi}(\cdot
  \mid y_i)} \log p_{\theta}(y_i \mid z_i)
\end{align*}
For this, we can use Monte Carlo averaging. First write $z_i \sim
N(\mu_{\phi}(y_i), \Sigma_{\phi}(y_i))$ as
\begin{align*}
  z_i = \mu_{\phi}(y_i) + \sigma_{\phi}(y_i) \odot u_i ~~
  \text{where}~ u_i \overset{\text{i.i.d}}{\sim} N(0, I_d). 
\end{align*}
Here $\odot$ stands for the Hadamard product (elementwise
multiplication). Note that the components of $\Sigma_{\phi}(y_i)^{1/2}
u_i$ are simply $\sigma_{\phi, j}(y_i) u_{i, j}$. Thus
\begin{align*}
   \E_{z_i \sim q_{\phi}(\cdot
  \mid y_i)} \log p_{\theta}(y_i \mid z_i) &=     \E_{u_i \sim N(0,
  I_d)} \log p_{\theta}(y_i \mid \mu_{\phi}(y_i) + \sigma_{\phi}(y_i)
                                             \odot u_i)  \\
  &\approx \frac{1}{L} \sum_{l=1}^L \log p_{\theta}(y_i \mid \mu_{\phi}(y_i) + \sigma_{\phi}(y_i)
                                             \odot u_i^{(l)}) 
\end{align*}
where $u_i^{(l)}, 1 \leq l \leq L, 1 \leq i \leq n$ are i.i.d samples
generated from $N(0, I_d)$.  Thus the ELBO becomes:
\begin{align*}
  & \text{ELBO}(\phi, \theta) \\
  &= \frac{1}{L} \sum_{i=1}^n \sum_{l = 1}^L \sum_{j=1}^{d_y}
    \left\{y_{i, j} \log \sigma\left( f_{\theta, j}(\mu_{\phi}(y_i) +
    \sigma_{\phi}(y_i) \odot u_i^{(l)})\right) \right. \\ & \left.+ (1 - y_{i, j}) \log \left[1 -
    \sigma\left( f_{\theta, j}(\mu_{\phi}(y_i) + 
    \sigma_{\phi}(y_i) \odot u_i^{(l)}) \right) 
    \right] \right\} \\
  &- \frac{1}{2} \sum_{i=1}^n \sum_{j=1}^d \left\{\sigma^2_{\phi,
    j}(y_i) + \mu^2_{\phi, j}(y_i) - \log \sigma^2_{\phi, j}(y_i) - 1
    \right\}. 
\end{align*}
Here $\sigma$ denotes the sigmoid function while $\sigma_{\phi}$
denotes the variance parameter of the variational distribution
$q_{\phi}$.

We simply maximize $\text{ELBO}(\phi, \theta)$ jointly over $\phi$ and
$\theta$ to estimate them. We will parametrize $\mu_{\phi}$ and
$\sigma_{\phi}$ using neural networks. For example,
\begin{align*}
  &  \mu_{\phi}(y) = W_{\mu} \text{ReLU}(Wy + w) + b_{\mu}  \\
  & \log \sigma^2_{\phi}(y) = W_{\sigma} \text{ReLU}(Wy + w) + b_{\sigma},
\end{align*}
with $\phi = (W_{\mu}, W, w, b_{\mu}, W_{\sigma}, b_{\sigma})$. Recall
also that $f_{\theta}$ will also be parametrized by a neural network
as in \eqref{fth_nn}.

With these parametrizations, $\text{ELBO}(\phi, \theta)$ can be
maximized using optimization software such as PyTorch. 



%\begin{thebibliography}{99}

%\bibitem[Neal and Hinton(1998)]{neal1998view}
%Neal, R.~M. and Hinton, G.~E. (1998).
%\newblock A view of the EM algorithm that justifies incremental, sparse, and other variants.
%\newblock In M.~I. Jordan (ed.), \emph{Learning in Graphical Models},
%pages 355--368. Springer Netherlands.  

%  \bibitem[Neal(2011)]{neal2011mcmc}
%Radford M. Neal.
%\newblock MCMC using Hamiltonian dynamics.
%\newblock In \emph{Handbook of Markov Chain Monte Carlo}, pages 47--95.
%\newblock Chapman and Hall/CRC, 2011.

%\bibitem[Girolami and Calderhead(2011)]{girolami2011riemann}
%Mark Girolami and Ben Calderhead.
%\newblock Riemann manifold Langevin and Hamiltonian Monte Carlo methods.
%\newblock \emph{Journal of the Royal Statistical Society: Series B (Statistical Methodology)}, 73(2):123--214, 2011.

%  \bibitem[Lai et~al.(2025)Lai, Song, Kim, Mitsufuji, and Ermon]{lai2025principles}
%Lai, C.-H., Song, Y., Kim, D., Mitsufuji, Y., and Ermon, S. (2025).
%\newblock The principles of diffusion models.
%\newblock \emph{arXiv preprint arXiv:2510.21890}.

%\bibitem[Glatt-Holtz et~al.(2024)]{glatt2024sacred}
%Nathan E. Glatt-Holtz, Andrew J. Holbrook, Justin A. Krometis, Cecilia F. Mondaini, and Ami Sheth.
%\newblock Sacred and Profane: from the Involutive Theory of MCMC to Helpful Hamiltonian Hacks.
%\newblock \emph{arXiv preprint arXiv:2410.17398}, 2024.

%\bibitem[Roberts and Rosenthal(2004)]{RobertsRosenthal2004}
%Roberts, G. O. and Rosenthal, J. S. (2004).
%General state space Markov chains and MCMC algorithms.
%\textit{Probability Surveys}, \textbf{1}, 20--71.

%\bibitem[Meyn and Tweedie(2009)]{MeynTweedie2009}
%Meyn, S. P. and Tweedie, R. L. (2009).
%\textit{Markov Chains and Stochastic Stability}, 2nd ed.
%Cambridge University Press.


  
%\bibitem[Chib and Greenberg(1995)]{chib1995}
%Chib, S. and Greenberg, E. (1995).
%Understanding the Metropolis--Hastings algorithm.
%\textit{The American Statistician}, \textbf{49}, 327--335.


%\bibitem[MacKay and Peto(1995)]{MacKay1995HierarchicalDirichlet}
%D.~J.~C. MacKay and L.~C. Peto,
%\newblock ``A hierarchical Dirichlet language model,''
%\newblock \emph{Natural Language Engineering}, vol.~1,
%no.~3, pp.~289--308, 1995.

%\bibitem[Rubin(1981)]{Rubin1981BayesianBootstrap}
%D.~B. Rubin,
%\newblock ``The Bayesian bootstrap,''
%\newblock \emph{The Annals of Statistics}, vol.~9,
%no.~1, pp.~130--134, 1981.
  
%\bibitem[Fisher(1934)]{Fisher1934}
%R.~A. Fisher,
%\newblock ``Probability, likelihood and quantity of information in the logic of
%uncertain inference,''
%\newblock \emph{Proceedings of the Royal Society of London. Series~A,
%Containing Papers of a Mathematical and Physical Character}, vol.~146,
%no.~856, pp.~1--8, 1934.

%\bibitem[Efron(2010)]{Efron}
%Efron, B. (2010)
%\textit{Large-scale inference: empirical Bayes methods for estimation,
%testing, and prediction}. Cambridge University Press, 2010.

%\bibitem[Shumway and Stoffer(2017)]{ShumwayStoffer}
%Shumway, R. H. and Stoffer, D. S. (2017)
%\textit{Time Series Analysis and Its Applications: With R
%  Examples}. Springer, 4th edition. 
  
%     \bibitem[Jaynes(2003)]{Jaynes} Jaynes, E. T, (2003)
%  \textit{Probability theory: the logic of science}. Cambridge
%  University Press, 2003.
 
%  \bibitem[Sinai(1992)]{Sinai} Sinai, Y. G, (1992)
%  \textit{Probability theory: an introductory course}. Springer
%  Verlag, 1992.  

%\bibitem[{deGroot(1986)}]{DeGrootBlackwell}DeGroot, M. H. (1986)
%       A conversation with David Blackwell.
%\newblock \emph{Statistical Science}, \textbf{1}, 40--53.

%         \bibitem[Mosteller(1987)]{Mosteller} Mosteller, F, (1987)
%  \textit{Fifty challenging problems in probability with
%    solutions}. Courier Corporation, 1987.

%    \bibitem[MacKay(2003)]{MacKay} MacKay, D. J. C., (2003)
%  \textit{Information theory, inference, and learning
%  algorithms}. Cambridge University Press, 2003.

%\bibitem[{Lindley(2013)}]{Lindley}Lindley, D. V. (2013)
%   \textit{Understanding Uncertainty}. John Wiley and sons, 2013.


%\end{thebibliography}







\end{document}

%%% Local Variables:
%%% mode: latex
%%% TeX-master: t
%%% End:
